\begin{abstract}
We characterize the attack surface of generative search engines (GSEs) to poisoning attacks in the political domain, from the perspectives of citation source selection and user personalization.
GSEs integrate web search and answer generation using large language models, playing a crucial role in how users access information, particularly in citizens' decision-making in politics.
Because anyone can publish content on the web, GSEs are vulnerable to poisoning attacks that manipulate citations to undermine reliable information delivery---a threat with direct implications for democracy.
Existing studies on citation evaluation focus on how faithfully answers reflect cited content, leaving two critical aspects unexamined: which web sources are selected as citations and how resilient those selections are to adversarial manipulation, and how user personalization affects citation behavior.
To fill this gap, we introduce an evaluation framework that characterizes the attack surface of GSEs to poisoning attacks.
Our contributions are twofold: (1) we propose a novel metric, \emph{content-injection barrier}, which quantifies the difficulty of injecting content into a source with a given level of publisher authority; and (2) we inject personas into system prompts to reveal how personalization shifts citation behavior.
We conduct experiments in the political domains of Japan and the United States using our framework.
Our results show that (a) for users with no prior knowledge of politics, GSEs tend to cite sources with low \emph{content-injection barriers}, making them susceptible to poisoning attacks; and (b) for users with prior knowledge, GSEs consistently cite sources with high \emph{content-injection barriers}, regardless of the model or political ideology.
\end{abstract}